%==========================
%          RESUMEN
%==========================
\chapter*{\hfill{\centering Resumen}\hfill}
\addcontentsline{toc}{chapter}{Resumen}  %Retirar comentario si quiere que aparesca en la lista del Indice general
\vspace{0.5cm}
\justify
\normalsize

El presente trabajo de investigación analiza analítica y perturbativamente la dinámica geodésica de fotones y la estructura de la sombra en el espaciotiempo de Kerr-Newman bajo la aproximación de carga eléctrica débil ($\delta \equiv \frac{Q^2}{4\pi\varepsilon_0 G M^2} \ll 1$), formulado rigurosamente en unidades del Sistema Internacional (SI) y empleando la convención de firma métrica de Minkowski $(+,-,-,-)$. Se deduce paso a paso el potencial radial efectivo $\mathcal{R}(r)$ a partir del formalismo de Hamilton-Jacobi y la función de acción $S$, obteniendo las constantes conservadas de energía $E$, momento angular azimutal $L_z$ y la constante de Carter $\mathcal{K}$. Definiendo los parámetros de impacto reducidos $\Phi \equiv \frac{c L_z}{E}$ y $\mathcal{Q} \equiv \frac{c^2 \mathcal{K}}{E^2}$, y estableciendo las condiciones de órbita fotónica esférica inestable $\mathcal{R}(r_{\mathrm{ph}}) = 0$ y $\left.\frac{d\mathcal{R}}{dr}\right|_{r_{\mathrm{ph}}} = 0$, se extiende la metodología de Kobialko para métricas rotacionales no diagonales a partir del parámetro crítico $b_c^2 = \mathcal{Q}_c + \Phi_c^2$. Se deriva la expansión perturbativa del radio efectivo equivalente $R_{\mathrm{eff}}(\delta)/R_{\mathrm{eff}}(0) = 1 - \frac{1}{6}\delta - \frac{7}{216}\delta^2$, el coeficiente de segundo orden $D = 1/27$ en $R_{\mathrm{eff}}^2$ que aísla la firma unívoca de electrovacío frente a modelos de gravedad modificada, y la razón de desplazamiento del primer nulo de visibilidad $u_1(\delta)/u_1(0) = 1 + \frac{1}{6}\delta + \frac{13}{216}\delta^2$. Asimismo, se demuestra que mientras en el límite estático de Schwarzschild el factor de adelgazamiento de Sneppen ($f_{\mathrm{thin}} = e^{-2\pi} \approx 1/535$) vuelve inobservable el subanillo $n=1$, en Kerr-Newman el momento angular degrada el exponente de Lyapunov ($\gamma_\phi < 1$), ensanchando orbitalmente el primer subanillo de fotones. Finalmente, se integran las geodésicas nulas mediante el algoritmo Runge-Kutta de cuarto orden (RK4) en Python y se contrastan los perfiles teóricos de visibilidad interferométrica $|V(u)|$ en el plano de Fourier con las observaciones del \textit{Event Horizon Telescope} (EHT) para $\text{Sgr A}^*$ (repositorio 2024-D02-01) y $\text{M87}^*$ (repositorio 2024-D01-01). Se encuentra que para $\text{Sgr A}^*$, el modelo de Kerr-Newman con $\delta \approx 0.290$ desplaza el primer nulo de visibilidad en un $+5.42\%$, reduciendo la bondad de ajuste en $\Delta\chi^2_\nu = -4.32$ ($\Delta\chi^2 = -40\,166$) frente a Schwarzschild; mientras que en $\text{M87}^*$, el modelo simétrico 1D empeora ($\Delta\chi^2_\nu = +10.74$) debido a la contaminación por la emisión extendida del chorro relativista. Se concluye que el ajuste en $\text{Sgr A}^*$ actúa como un compensador de escala geométrica en el modelo simétrico 1D y no constituye una detección física directa de carga, dada la dominancia de las incertidumbres sistemáticas actuales en la razón masa-distancia ($\sim 11.8\%$).

%=================================
% Palabras claves
%=================================
\vspace{0.5cm}
\raggedright 
{\bf \textit{Palabras clave}:} Agujero negro de Kerr-Newman, Carga débil, Esfera de fotones, Radio efectivo de la sombra, Constante de Carter, Interferometría VLBI, Event Horizon Telescope, Exponente de Lyapunov.
